\section{Ethics and Dual Use}
\label{sec:ethics}

\paragraph{What this system can and cannot do.} The detector recognises visible
mine-like objects and surface indicators in downward-looking colour imagery. It
cannot detect a buried object, because the sensor cannot see one. It has not been
flown, field-tested, or validated against ground truth collected in a real
clearance operation. Nothing in this paper supports a statement about detection
probability in the field, and we do not make one. The correct description is a
computer-vision benchmark on a public, publicly derived image corpus.

\paragraph{Dual use.} The same capability that helps a clearance team prioritise
ground could help someone avoid marked hazards. We publish standard detection
metrics, an audit methodology, and trained weights; we publish no targeting
guidance, no geographic coordinates, and no operational deployment procedure.
The dual-use risk is real but generic to visible-object detection, and we judge
that the methodological contribution---how to audit a leaky split and why to
report operating points---outweighs it. We note that the weights are released
under AGPL-3.0, which requires derivative services to publish their source; this
is a licence choice, not an ethical safeguard.

\paragraph{The cost of the miss rate.} At the primary operating point the
two-stage arm recalls $0.964$ of annotated instances, which on this test set means
nine missed objects out of $253$. Those nine are analysed in
Table~\ref{tab:failures}: five are camouflaged or low-contrast, two are small
objects, and two are the objectively wrong labels removed under D023. A
pre-correction recall of $0.964$ reads as a good number in a table; after the
two label corrections, the corresponding value is $0.972$. In a clearance context it
means that roughly one in twenty-eight annotated objects is not flagged, and the
misses are concentrated in exactly the low-contrast cases where a human reviewer
is most likely to defer to the model. No metric in this paper should be read as
saying the system is safe to rely on.

\paragraph{Data ethics.} The corpus is derived from public dataset exports and
contains no human subjects, no personal data, and no imagery of identified
individuals. We did not collect new data. The original licences were not recorded
at download time; we therefore redistribute no images (Section~\ref{sec:data}).
Where figures include source imagery, they are reproduced for scholarly comment
at low resolution and are not a dataset release.

\paragraph{Reporting.} We report the results that weaken our own claims---the
seed-2 failure of the two-stage pipeline, the negative leakage result, the
degenerate image-level baseline, the tuned single-stage threshold that beats the
pipeline---in the main text rather than in an appendix. The reason is practical:
a reader who is deciding whether to trust the positive results needs to know how
the negative ones were handled.
